\section{Muebles propileno}

\subsection{Presentación del caso}
Un reciclador de polipropileno (PP) dispone semanalmente de 900 kg de residuo de PP con los que puede obtener dos 
productos diferentes de mobiliario urbano: bancos y mesas. Su principal limitación de fabricación está en las 80 horas 
máquina semanales disponibles para la fabricación. 

Conocidos los requisitos unitarios de los dos posibles productos a fabricar (en horas-máquina de fabricación y en kg de 
residuo de PP), así como sus contribuciones unitarias al beneficio en euros, se plantea el siguiente modelo de 
programación lineal para determinar el plan semanal óptimo de producción.  

\begin{center}
\begin{tabular}{l c c c}
\hline
Producto & h·m/unidad & kg/unidad & €/unidad \\
\hline
Banco & 2 & 30 & 300 \\
Mesa  & 2 & 15 & 400 \\
\hline
\end{tabular}
\end{center}

Además, existe el compromiso de suministrar al menos 10 bancos semanales y el mercado de mesas está limitado a 25 
unidades por semana .

\subsection{Formulación explícita}

\paragraph{Variables}\mbox{}\\
\begin{tabular}{l c l}
    $x_1$  & : & número de bancos producidos semanalmente\\
    $x_2$  & : & número de mesas producidas semanalmente\\
\end{tabular}

\paragraph{Restricciones}\mbox{}\\

Disponibilidad máxima de horas máquina:
\begin{equation}
    2x_1 + 2x_2 \leq 80
\end{equation}

Disponibilidad máxima de polipropileno:
\begin{equation}
    30x_1 + 15x_2 \leq 900
\end{equation}

Compromiso mínimo de producción de bancos:
\begin{equation}
    x_1 \geq 10
\end{equation}

Limitación máxima de mercado de mesas:
\begin{equation}
    x_2 \leq 25
\end{equation}

Restricciones de no negatividad:
\begin{equation}
    x_1, x_2 \geq 0
\end{equation}

\paragraph{Función objetivo}\mbox{}\\
Beneficio total semanal:
\begin{equation}
    \mbox{max. } z = 300x_1 + 400x_2
\end{equation}

\paragraph{Modelo completo}\mbox{}\\
\begin{align}
\mbox{max.} \quad & 300x_1 + 400x_2 \notag \\
\mbox{s. a.:} \quad & 2x_1 + 2x_2 \leq 80 \notag \\
             & 30x_1 + 15x_2 \leq 900 \notag \\
             & x_1 \geq 10 \notag \\
             & x_2 \leq 25 \notag \\
             & x_1, x_2 \geq 0
\end{align}

\subsection{Formulación generalizada}

\paragraph{Conjuntos}\mbox{}\\
\begin{tabular}{l c l}
    $\mathcal{P}$  & : & tipos de productos de mobiliario urbano (bancos, mesas) \\
\end{tabular}

\paragraph{Parámetros}\mbox{}\\
\begin{tabular}{l c l}
    $H$  & : & número máximo de horas máquina disponibles por semana \\
    $W_p$  & : & horas máquina requeridas para producir una unidad de producto $p \in \mathcal{P}$ \\
    $K$  & : & cantidad máxima de polipropileno disponible semanalmente (kg) \\
    $M_p$  & : & kg de polipropileno requeridos para producir una unidad de producto $p \in \mathcal{P}$ \\
    $B_p$  & : & beneficio unitario por producir una unidad de producto $p \in \mathcal{P}$ \\
    $L_b$  & : & compromiso mínimo de producción de bancos \\
    $U_m$  & : & máximo de mesas que el mercado puede absorber semanalmente \\
\end{tabular}

\paragraph{Variables}\mbox{}\\
\begin{tabular}{l c l}
    $x_p$  & : & número de unidades producidas de producto $p \in \mathcal{P}$ \\
\end{tabular}

\paragraph{Restricciones}\mbox{}\\

Limitación de horas máquina:
\begin{equation}
    \sum_{p \in \mathcal{P}} W_p x_p \leq H
\end{equation}

Limitación de polipropileno:
\begin{equation}
    \sum_{p \in \mathcal{P}} M_p x_p \leq K
\end{equation}

Compromiso de bancos:
\begin{equation}
    x_{\text{banco}} \geq L_b
\end{equation}

Limitación de mercado de mesas:
\begin{equation}
    x_{\text{mesa}} \leq U_m
\end{equation}

No negatividad:
\begin{equation}
    x_p \geq 0 \quad \forall p \in \mathcal{P}
\end{equation}

\paragraph{Función objetivo}\mbox{}\\
\begin{equation}
    \mbox{max. } z = \sum_{p \in \mathcal{P}} B_p x_p
\end{equation}

\paragraph{Modelo completo}\mbox{}\\
\begin{align}
\mbox{max.} \quad & \sum_{p \in \mathcal{P}} B_p x_p \notag \\
\mbox{s. a.:} \quad & \sum_{p \in \mathcal{P}} W_p x_p \leq H \notag \\
             & \sum_{p \in \mathcal{P}} M_p x_p \leq K \notag \\
             & x_{\text{banco}} \geq L_b \notag \\
             & x_{\text{mesa}} \leq U_m \notag \\
             & x_p \geq 0 \quad \forall p \in \mathcal{P} \notag
\end{align}
